\section{Limitations}
\label{sec:limits}

\num{plantvillage:seeds} class partitions are few. The spread over partitions in the tables mixes how hard a partition is with training noise, and it is itself estimated from five numbers. Every hyper-parameter, including $\lambda$, was fixed in advance and shared by all scores, and the only value we tuned (the episodic learning rate, on validation classes) was tuned on the dense network. A score that needs its own learning rate or recovery schedule did not get its best case.

PlantVillage images are leaves on a plain background, and the dataset holds several photographs of the same leaf. We did not use leaf identities to keep support and query images apart, so near-duplicates can raise the PlantVillage numbers; the PlantDoc and EuroSAT-MS results do not depend on that. PlantDoc here is a cropped-leaf release with one dropped class, and its concepts overlap those of PlantVillage. EuroSAT-MS has ten classes and $64{\times}64$ patches, which allows a $5/5$ class split and nothing larger. No model uses coordinates, acquisition dates or any sensor beyond the thirteen optical bands, and none was tested on full scenes.

Cost is reported as parameters, multiply-accumulates, and latency on one server CPU and one GPU. We did not run on a Raspberry Pi, a Jetson or an accelerator on board a satellite, and we measured no energy, so we say nothing about those devices. Only ResNet-18 and only the inner convolution of each residual block were pruned. Quantisation, distillation, architecture search and recent few-shot learners were not compared; the claims are limited to channel scores and training schedules inside this one design.
